\input{../../preamble.tex}

\begin{document}
    
% Indicate your name and your student number here (e.g., A01xxx)
\header{Lab 1}{Wira Azmoon Ahmad}{A0149286R}

%======================================
\section{DSDA Model}

A* Search is used for this model, as the C-space is generally small. 
On the other hand, the state space can be large especially for the larger
COM1 map, so dynamic programming will be largely held back and not practical.
As such A* Search is used. For this Lab, the Manhattan Distance heuristic
is used as it is an admissible heuristic for discrete states. It is also
better than the Euclidean distance heuristic simply via the triangle inequality.

\section{CSDA Model}

The state space was discretized in both location and rotation.
For the $(x,y)$ coordinates, this was split into grid of sizes smaller
than 1. In the generated model the chosen value was half.
As for $\theta$, this was split into 8 to allow for finer control 
of the robot turning.
In the code, the discretization can be seen from the variables
$\texttt{NUM\_SPLIT\_V, NUM\_SPLIT\_W}$.
Hybrid A* Search is used with a combination of holonomic and
non-holonomic heuristics. 
Holonomic calculated the shortest path to the goal using BFS,
taking into account diagonal movement,
while non-holonomic used the Euclidean distance to the goal.
Both are easily admissible, and their maximum was used.

\section{DSPA Model}

The reward function used was closely modeled as
\[R(s,a)=\sum_{s'}T(s,a,s')R(s')\]

\[R(s) = \frac{o(s)R_{MAX}}{d(s) + 1}\]
\[o(s) =
    \begin{cases}
        1, & \text{if no obstacles in or around s}\\
        -1, & \text{if s is an obstacle itself}\\
        0, & \text{if an obstacle is a step away from s}
    \end{cases}
\]
where $d(s)$ is the distance to the goal determined by BFS.\\
It is chosen that $R_{MAX}=1$. 
For some goals, the $R(s), o(s)$ were modified to be
more or less risky, and empirically tested to determine the better settings.
The choice of the reward function largely rewards positions 
closer to the goal, and gives 0 reward for positions near obstacles.
The MDP is solved using Value Iteration, with tweaking of values such
as inflation and $\epsilon$ depending on the map.
This was chosen due to the large state space, while the dimension
of the C-space is small. Other MDP solutions would suffer from
dealing with a large state space.

\section{Lessons and Insights}

I found that even simple path planning was no easy feat, working through
the various solutions for each problem. There is a lot of hacking and
testing to see which algorithms and parameters fit which problem.
There is no one-size-fits-all solution, even using a similar model.

\end{document}